\section{Fusion Benefit vs.\ Hand Self-Occlusion}
\label{app:occlusion}

\begin{figure}[t]
  \centering
  \includegraphics[width=\linewidth]{figures/occlusion_fusion_gain.pdf}
  \caption{%
    Fusion benefit as a function of hand self-occlusion on the EgoEMG test set
    (17{,}168 samples, stride reduced to $\frac{1}{4}$ of the evaluation window for denser coverage).
    \textbf{(a)}~Mean joint-angle MAE of vision-only and EMG+Vision fusion across
    fixed-width (0.05) occlusion bins.
    The annotated $\Delta$ values indicate the fusion gain at the lowest and highest
    occlusion bins.
    \textbf{(b)}~Per-bin fusion gain ($\Delta = \text{MAE}_{\text{vision}} - \text{MAE}_{\text{fusion}}$)
    with LOESS trend line (Spearman $\rho = 0.049$, $p = 0.001$).
    Error bands in~(a) and the trend confidence in~(b) denote $\pm$1 SEM.
  }
  \label{fig:occlusion_gain}
\end{figure}

To understand when EMG--vision fusion provides the greatest benefit, we
stratify test-set predictions by the degree of hand self-occlusion.

\paragraph{Self-occlusion score.}
We quantify per-frame hand self-occlusion via z-buffer rasterization of the
ground-truth MANO mesh in camera space.  Concretely, for each test frame:
\begin{enumerate}[nosep]
  \item The MANO mesh vertices are transformed from world coordinates to camera
        coordinates using the calibrated extrinsic matrix ($\mathbf{v}_{\text{cam}} =
        \mathbf{R}_{C \leftarrow W}\,\mathbf{v}_W + \mathbf{t}_{C \leftarrow W}$).
  \item All mesh triangles are rasterized into a depth buffer (z-buffer) at the
        native video resolution via pinhole projection with the calibrated
        intrinsic matrix~$\mathbf{K}$.  For each pixel covered by a triangle, the
        interpolated depth is stored if it is smaller than the current buffer value.
  \item A vertex is deemed \emph{visible} if, within a $5 \times 5$ pixel
        neighbourhood of its projected location, any z-buffer entry agrees with
        the vertex depth to within $\epsilon = 5$\,mm.
  \item Each triangle's 3-D surface area (half the cross-product norm of two
        edge vectors) is distributed equally to its three vertices, yielding a
        per-vertex area weight~$w_i$.
  \item The \textbf{self-occlusion score} is defined as
        \begin{equation}
          s_{\text{occ}} = 1 - \frac{\sum_{i \in \mathcal{V}_{\text{vis}}} w_i}
                                     {\sum_{i=1}^{V} w_i}\,,
        \end{equation}
        where $\mathcal{V}_{\text{vis}}$ is the set of visible vertices.
        A score of~0 means the hand is fully visible; a score of~1 means it is
        fully occluded from the camera viewpoint.
\end{enumerate}

\paragraph{Results.}
Figure~\ref{fig:occlusion_gain}(a) shows that both vision-only and fusion MAE
increase with occlusion, but the gap between them---the fusion gain---grows
monotonically from $\Delta = 0.24\degree$ (occlusion $\approx 0.38$) to
$\Delta = 0.78\degree$ (occlusion $\approx 0.72$), confirming that EMG's
proprioceptive signal becomes increasingly valuable as visual evidence degrades.
At extreme occlusion ($> 0.80$), both modalities deteriorate and the gain
narrows slightly, suggesting a regime where neither modality alone is reliable.
Panel~(b) visualises the same trend at the per-bin level with a LOESS smooth,
showing a statistically significant positive correlation between occlusion and
fusion gain (Spearman $\rho = 0.049$, $p = 0.001$).
